\chapter{Value Functions and the Bellman Equations}

The return $G_t$ scores a trajectory. The interaction is an MDP: the policy $\pi$ selects the action; the environment replies with $(s',r)$ from $P$. That return is one number for one realized sequence. Take the taxi introduced in Chapter 1: from a junction it still faces many possible futures. The policy $\pi$ may leave more than one turn possible, and $P$ may send it to different streets with different fares. Two trips that start at the same junction can therefore finish with different returns: $G_t$ from that junction is a random variable. One realized return does not tell how good the junction is.

To say how good that junction is under a fixed $\pi$, take the expected return from there: the \textbf{state-value} $V^\pi$. To say how good a particular turn is, take the expected return after that turn: the \textbf{action-value} $Q^\pi$. Both are rewritten as \textbf{Bellman equations} (\citetxt{bellman1957}{Bellman, 1957}) that use only the immediate reward and the value at the next state. Averaging over $\pi$ \textbf{evaluates} that policy. Replacing the average by $\max_a$ is optimality: the rewrite that \textbf{improves} $\pi$.

The expected return from $s$ when actions afterward follow $\pi$ is the \textbf{state-value} $V^\pi(s)$.

\begin{equation}
V^\pi(s) = \mathbb{E}_\pi [G_t \mid S_t = s].
\tag{2.1}
\end{equation}

We also care how good a particular turn $a$ is at that junction. Fix $A_t=a$, then follow $\pi$ from the next state on: the expected return is the \textbf{action-value} $Q^\pi(s,a)$.

\begin{equation}
Q^\pi(s,a) = \mathbb{E}_\pi [G_t \mid S_t = s, A_t = a].
\tag{2.2}
\end{equation}

\section{Bellman expectation equations}
$V^\pi(s)$ is an expected return, and the return is an infinite discounted sum of future rewards. Computing that expectation by listing every possible future under $\pi$ and $P$ is not practical. Using $G_t=R_{t+1}+\gamma G_{t+1}$ in the definition of $V^\pi$, the Bellman expectation equation is obtained as follows.

\[
\begin{aligned}
V^\pi(s)
&= \mathbb{E}_\pi [G_t \mid S_t = s] \\
&= \mathbb{E}_\pi [R_{t+1} + \gamma G_{t+1} \mid S_t = s] \\
&= \mathbb{E}_\pi [R_{t+1} \mid S_t = s] + \gamma \mathbb{E}_\pi [G_{t+1} \mid S_t = s].
\end{aligned}
\]

So

\begin{equation}
V^\pi(s)
=
\mathbb{E}_\pi [R_{t+1} \mid S_t = s]
+
\gamma\, \mathbb{E}_\pi [G_{t+1} \mid S_t = s].
\tag{2.3}
\end{equation}

The first term is the expected immediate reward from $s$ under $\pi$. The second is the expected value of the next state, again under $\pi$. Expanding both over $\pi(a\mid s)$ and $P(s',r\mid s,a)$ is in \appref{B}, \eqref{eq:B.1}. The result is

\begin{equation}
V^\pi(s) = \sum_{a \in \mathcal{A}} \pi(a \mid s) \sum_{s', r} P(s',r \mid s,a) \left[ r + \gamma V^\pi(s') \right].
\tag{2.4}
\end{equation}

The inner piece $r+\gamma V^\pi(s')$ looks only one step ahead: $r$ is the immediate reward on the transition to $s'$ (a value of $R_{t+1}$), and $V^\pi(s')$ is the value at the next state. That one-step rewrite is a \textbf{backup}.

Equation (2.4) nests two averages---over actions under $\pi$, then over transitions under $P$. These are the same as $V^\pi$ through $Q^\pi$ (\appref{B}, \eqref{eq:B.14} and \eqref{eq:B.17}):

\[
V^\pi(s)
=
\sum_{a \in \mathcal{A}}
\pi(a \mid s)\, Q^\pi(s,a),
\qquad
Q^\pi(s,a)
=
\sum_{s',r}
P(s',r \mid s,a)
\bigl[
r + \gamma V^\pi(s')
\bigr].
\]

So the outer sum in (2.4) averages over actions under $\pi$, and the inner sum averages over the next state and reward under $P$.

Starting from $Q^\pi(s,a)$ and eliminating $V^\pi$ at the next state (\appref{B}, \eqref{eq:B.2}) gives

\begin{equation}
Q^\pi(s,a) = \sum_{s',r} P(s',r \mid s,a) \left[ r + \gamma \sum_{a'} \pi(a' \mid s') Q^\pi(s',a') \right].
\tag{2.5}
\end{equation}

\subsection{Matrix form}
The policy-averaged form of $V^\pi$ (\appref{B}, \eqref{eq:B.3}) is

\[
V^\pi(s)
=
R^\pi(s)
+
\gamma
\sum_{s'}
P^\pi(s' \mid s)\, V^\pi(s'),
\]

with $R^\pi$ and $P^\pi$ as in Section 1.3. For a finite state space of size $N$, stack $V^\pi(s)$ into a vector $\mathbf{V}^\pi\in\mathbb{R}^N$, stack $R^\pi(s)$ into $\mathbf{R}^\pi$, and let $\mathbf{P}^\pi$ be the $N\times N$ matrix with entry $(s,s')$ equal to $P^\pi(s'\mid s)$. The display above for every $s$ is the single matrix equation

\[
\mathbf{V}^\pi = \mathbf{R}^\pi + \gamma \mathbf{P}^\pi \mathbf{V}^\pi.
\]

Solve for $\mathbf{V}^\pi$:

\begin{equation}
\mathbf{V}^\pi - \gamma \mathbf{P}^\pi \mathbf{V}^\pi = \mathbf{R}^\pi
\implies
(I - \gamma \mathbf{P}^\pi)\mathbf{V}^\pi = \mathbf{R}^\pi
\implies
\mathbf{V}^\pi = (I - \gamma \mathbf{P}^\pi)^{-1} \mathbf{R}^\pi.
\tag{2.6}
\end{equation}

For $\gamma\in[0,1)$, $I-\gamma\mathbf{P}^\pi$ is invertible, so evaluating $\pi$ is one linear solve---not an iteration.

\begin{asidenote}[Note]
A two-junction taxi makes the arrays concrete. The states are the depot $D$ and the street $S$. Under $\pi$ the taxi always leaves the depot for the street (immediate reward $0$). From the street it collects a fare of $1$ and returns to the depot with probability $1/2$, or stays on the street with reward $0$ with probability $1/2$. Then

\[
\mathbf{P}^\pi
=
\begin{bmatrix}
0 & 1 \\
\tfrac{1}{2} & \tfrac{1}{2}
\end{bmatrix},
\qquad
\mathbf{R}^\pi
=
\begin{bmatrix}
0 \\
\tfrac{1}{2}
\end{bmatrix},
\]

with rows and columns ordered $(D,S)$. Take $\gamma=1/2$. The two scalar copies of (2.4) are

\[
\begin{aligned}
V^\pi(D)
&=
0 + \tfrac{1}{2} V^\pi(S), \\
V^\pi(S)
&=
\tfrac{1}{2} + \tfrac{1}{2}
\bigl(
\tfrac{1}{2} V^\pi(D) + \tfrac{1}{2} V^\pi(S)
\bigr).
\end{aligned}
\]

The same system is (2.6):

\[
I - \gamma \mathbf{P}^\pi
=
\begin{bmatrix}
1 & -\tfrac{1}{2} \\
-\tfrac{1}{4} & \tfrac{3}{4}
\end{bmatrix},
\qquad
\mathbf{V}^\pi
=
(I - \gamma \mathbf{P}^\pi)^{-1} \mathbf{R}^\pi
=
\begin{bmatrix}
\tfrac{2}{5} \\
\tfrac{4}{5}
\end{bmatrix}.
\]

So $V^\pi(D)=2/5$ and $V^\pi(S)=4/5$. Substituting back into $\mathbf{R}^\pi+\gamma\mathbf{P}^\pi\mathbf{V}^\pi$ recovers the same vector.
\end{asidenote}

\section{Bellman optimality equations}
The Bellman expectation equations are for one fixed $\pi$. We now want values that are best in every state: the \textbf{optimal} state-value and action-value

\begin{equation}
V^*(s) = \max_\pi V^\pi(s),
\qquad
Q^*(s,a) = \max_\pi Q^\pi(s,a).
\tag{2.7}
\end{equation}

Under an optimal policy the agent no longer averages actions with $\pi(a\mid s)$; it takes an action that attains the maximum. The same one-step backup as (2.4), with that $\max$ in place of the average over $\pi$, is the \textbf{Bellman optimality equation} (\appref{C}, \eqref{eq:C.2}):

\begin{equation}
V^*(s) = \max_{a \in \mathcal{A}} \sum_{s',r} P(s',r \mid s,a) \left[ r + \gamma V^*(s') \right].
\tag{2.8}
\end{equation}

The $\max$ means: score every action by that sum and keep the best. That is only possible once the next-state values $V^*(s')$ are known; those values depend on other states the same way, so (2.8) couples the whole table and does not yield $V^*$ by one pass of ``try every $a$.''

\begin{equation}
Q^*(s,a) = \sum_{s',r} P(s',r \mid s,a) \left[ r + \gamma \max_{a'} Q^*(s',a') \right].
\tag{2.9}
\end{equation}

From (2.8), an optimal policy is greedy on those one-step scores:

\begin{equation}
\pi^*(s)
=
\arg\max_{a \in \mathcal{A}}
\sum_{s',r}
P(s',r \mid s,a)
\left[
r + \gamma V^*(s')
\right].
\tag{2.10}
\end{equation}

That choice uses $P$ even after $V^*$ is known. Computing $Q^*$ from (2.9) also uses $P$. Once the table $Q^*$ is known, the same policy is an $\arg\max$ over $a$ and does not use $P$:

\begin{equation}
\pi^*(s) = \arg\max_{a \in \mathcal{A}} Q^*(s,a).
\tag{2.11}
\end{equation}

The $\arg\max$ can contain more than one action: two turns with the same score. A greedy $\pi^*$ always takes one of them. Randomizing between those two is also optimal. Taking a turn whose score is strictly smaller is not (\appref{C}, \hyperref[app:C.4]{C.4}).

Equations (2.4)--(2.6) are linear in $V^\pi$, so a fixed policy is one matrix solve; (2.8)--(2.9) contain a $\max$, so that solve no longer applies and $V^*$ needs iteration (Chapter 3).
